%%%%%%%%%%%%%%%%%%%%%%%%%%%%%%%%%%%%%%%%%%%%%%%%%%%%%%%%%%%%
% 9.2 TRANSFORM METHODS
% The Composition of Advanced Mathematics
%%%%%%%%%%%%%%%%%%%%%%%%%%%%%%%%%%%%%%%%%%%%%%%%%%%%%%%%%%%%

\section{Transform Methods}
\label{sec:transform-methods}

\prerequisite{
Ordinary differential equations, improper integrals, exponential and
trigonometric functions, complex numbers, sequences and series, linear
algebra, and basic piecewise-defined functions.
}

\learningobjectives{
By the end of this section, the reader should be able to:
\begin{itemize}
    \item explain the purpose of integral transforms;
    \item define the Laplace transform;
    \item determine conditions under which a Laplace transform exists;
    \item compute elementary Laplace transforms;
    \item use linearity of the Laplace transform;
    \item use transform tables effectively;
    \item derive and use first- and second-shifting theorems;
    \item transform derivatives and integrals;
    \item solve linear initial-value problems using Laplace transforms;
    \item handle polynomial, exponential, and trigonometric forcing;
    \item understand Heaviside step functions;
    \item represent piecewise forcing using step functions;
    \item transform delayed functions;
    \item understand the Dirac delta distribution conceptually;
    \item solve impulsively forced differential equations;
    \item define convolution;
    \item use the convolution theorem;
    \item invert rational Laplace transforms using partial fractions;
    \item solve systems of differential equations using transforms;
    \item understand transfer functions;
    \item understand impulse and step responses;
    \item recognize poles and their relationship to system behavior;
    \item understand the bilateral Laplace transform conceptually;
    \item define the Fourier transform under a stated convention;
    \item understand the relation between Laplace and Fourier transforms;
    \item recognize sine and cosine transforms;
    \item understand transform methods for differential equations;
    \item recognize the \(z\)-transform as a discrete analogue;
    \item connect transform methods with ODEs, PDEs, signal processing,
          control theory, and numerical mathematics.
\end{itemize}
}

%%%%%%%%%%%%%%%%%%%%%%%%%%%%%%%%%%%%%%%%%%%%%%%%%%%%%%%%%%%%
\subsection{Why Use Transform Methods?}
\label{subsec:why-transform-methods}
%%%%%%%%%%%%%%%%%%%%%%%%%%%%%%%%%%%%%%%%%%%%%%%%%%%%%%%%%%%%

A transform converts a function from one representation into another.

The main idea is

\[
\boxed{
\text{difficult operation in the original domain}
\longrightarrow
\text{simpler operation in the transform domain}.
}
\]

For differential equations, derivatives can often be transformed into
algebraic expressions.

Thus

\[
\boxed{
\text{differential equation}
\longrightarrow
\text{algebraic equation}
\longrightarrow
\text{inverse transform}.
}
\]

%%%%%%%%%%%%%%%%%%%%%%%%%%%%%%%%%%%%%%%%%%%%%%%%%%%%%%%%%%%%
\subsection{Integral Transforms}
\label{subsec:integral-transforms}
%%%%%%%%%%%%%%%%%%%%%%%%%%%%%%%%%%%%%%%%%%%%%%%%%%%%%%%%%%%%

A general integral transform has form

\[
\boxed{
F(s)
=
\int
K(s,t)f(t)\,dt,
}
\]

where

\[
K(s,t)
\]

is called the kernel of the transform.

Different kernels produce different transforms.

Important examples include:

\[
\boxed{
\text{Laplace transform},
}
\]

\[
\boxed{
\text{Fourier transform},
}
\]

and

\[
\boxed{
\text{sine and cosine transforms}.
}
\]

%%%%%%%%%%%%%%%%%%%%%%%%%%%%%%%%%%%%%%%%%%%%%%%%%%%%%%%%%%%%
\subsection{Laplace Transform}
\label{subsec:laplace-transform}
%%%%%%%%%%%%%%%%%%%%%%%%%%%%%%%%%%%%%%%%%%%%%%%%%%%%%%%%%%%%

\begin{definitionbox}{Laplace Transform}
For a function \(f(t)\) defined for \(t\geq0\), its Laplace transform is

\[
\boxed{
\mathcal{L}\{f(t)\}(s)
=
F(s)
=
\int_0^\infty
e^{-st}f(t)\,dt,
}
\]

whenever the improper integral converges.
\end{definitionbox}

The variable \(t\) is usually the original time variable.

The variable \(s\) is the transform variable.

%%%%%%%%%%%%%%%%%%%%%%%%%%%%%%%%%%%%%%%%%%%%%%%%%%%%%%%%%%%%
\subsection{Transform Notation}
\label{subsec:laplace-transform-notation}
%%%%%%%%%%%%%%%%%%%%%%%%%%%%%%%%%%%%%%%%%%%%%%%%%%%%%%%%%%%%

Several equivalent notations are common:

\[
\boxed{
\mathcal{L}\{f(t)\}
=
F(s),
}
\]

or simply

\[
\boxed{
\mathcal{L}\{f\}
=
F.
}
\]

The inverse transform is written

\[
\boxed{
f(t)
=
\mathcal{L}^{-1}\{F(s)\}.
}
\]

%%%%%%%%%%%%%%%%%%%%%%%%%%%%%%%%%%%%%%%%%%%%%%%%%%%%%%%%%%%%
\subsection{Existence of the Laplace Transform}
\label{subsec:laplace-existence}
%%%%%%%%%%%%%%%%%%%%%%%%%%%%%%%%%%%%%%%%%%%%%%%%%%%%%%%%%%%%

A standard sufficient condition is that \(f\) be piecewise continuous on
every finite interval and of exponential order.

%%%%%%%%%%%%%%%%%%%%%%%%%%%%%%%%%%%%%%%%%%%%%%%%%%%%%%%%%%%%
\subsection{Exponential Order}
\label{subsec:exponential-order}
%%%%%%%%%%%%%%%%%%%%%%%%%%%%%%%%%%%%%%%%%%%%%%%%%%%%%%%%%%%%

\begin{definitionbox}{Exponential Order}
A function \(f(t)\) is of exponential order \(a\) if there exist
constants \(M>0\) and \(T\geq0\) such that

\[
\boxed{
|f(t)|
\leq
Me^{at}
}
\]

for all \(t\geq T\).
\end{definitionbox}

If \(f\) is piecewise continuous and of exponential order \(a\), then its
Laplace transform exists for sufficiently large

\[
\boxed{
s>a.
}
\]

%%%%%%%%%%%%%%%%%%%%%%%%%%%%%%%%%%%%%%%%%%%%%%%%%%%%%%%%%%%%
\subsection{Linearity}
\label{subsec:laplace-linearity}
%%%%%%%%%%%%%%%%%%%%%%%%%%%%%%%%%%%%%%%%%%%%%%%%%%%%%%%%%%%%

The Laplace transform is linear:

\[
\boxed{
\mathcal{L}
\{
af+bg
\}
=
a\mathcal{L}\{f\}
+
b\mathcal{L}\{g\}.
}
\]

Thus if

\[
\mathcal{L}\{f\}=F
\]

and

\[
\mathcal{L}\{g\}=G,
\]

then

\[
\boxed{
\mathcal{L}
\{
af+bg
\}
=
aF+bG.
}
\]

%%%%%%%%%%%%%%%%%%%%%%%%%%%%%%%%%%%%%%%%%%%%%%%%%%%%%%%%%%%%
\subsection{Transform of a Constant}
\label{subsec:laplace-constant}
%%%%%%%%%%%%%%%%%%%%%%%%%%%%%%%%%%%%%%%%%%%%%%%%%%%%%%%%%%%%

For

\[
f(t)=1,
\]

\[
\mathcal{L}\{1\}
=
\int_0^\infty
e^{-st}\,dt.
\]

For \(s>0\),

\[
\boxed{
\mathcal{L}\{1\}
=
\frac1s.
}
\]

%%%%%%%%%%%%%%%%%%%%%%%%%%%%%%%%%%%%%%%%%%%%%%%%%%%%%%%%%%%%
\subsection{Transform of a Power}
\label{subsec:laplace-power}
%%%%%%%%%%%%%%%%%%%%%%%%%%%%%%%%%%%%%%%%%%%%%%%%%%%%%%%%%%%%

For nonnegative integer \(n\),

\[
\boxed{
\mathcal{L}\{t^n\}
=
\frac{n!}{s^{n+1}}.
}
\]

More generally, for \(\alpha>-1\),

\[
\boxed{
\mathcal{L}\{t^\alpha\}
=
\frac{
\Gamma(\alpha+1)
}{
s^{\alpha+1}
}.
}
\]

%%%%%%%%%%%%%%%%%%%%%%%%%%%%%%%%%%%%%%%%%%%%%%%%%%%%%%%%%%%%
\subsection{Transform of an Exponential}
\label{subsec:laplace-exponential}
%%%%%%%%%%%%%%%%%%%%%%%%%%%%%%%%%%%%%%%%%%%%%%%%%%%%%%%%%%%%

For

\[
f(t)=e^{at},
\]

\[
\mathcal{L}\{e^{at}\}
=
\int_0^\infty
e^{-(s-a)t}\,dt.
\]

Therefore,

\[
\boxed{
\mathcal{L}\{e^{at}\}
=
\frac1{s-a},
\qquad
s>a.
}
\]

%%%%%%%%%%%%%%%%%%%%%%%%%%%%%%%%%%%%%%%%%%%%%%%%%%%%%%%%%%%%
\subsection{Transform of Sine and Cosine}
\label{subsec:laplace-trigonometric}
%%%%%%%%%%%%%%%%%%%%%%%%%%%%%%%%%%%%%%%%%%%%%%%%%%%%%%%%%%%%

The standard transforms are

\[
\boxed{
\mathcal{L}\{\cos bt\}
=
\frac{s}{
s^2+b^2
},
}
\]

and

\[
\boxed{
\mathcal{L}\{\sin bt\}
=
\frac{b}{
s^2+b^2
}.
}
\]

%%%%%%%%%%%%%%%%%%%%%%%%%%%%%%%%%%%%%%%%%%%%%%%%%%%%%%%%%%%%
\subsection{Transform of Hyperbolic Functions}
\label{subsec:laplace-hyperbolic}
%%%%%%%%%%%%%%%%%%%%%%%%%%%%%%%%%%%%%%%%%%%%%%%%%%%%%%%%%%%%

Using exponential definitions,

\[
\boxed{
\mathcal{L}\{\cosh at\}
=
\frac{s}{
s^2-a^2
},
}
\]

and

\[
\boxed{
\mathcal{L}\{\sinh at\}
=
\frac{a}{
s^2-a^2
}.
}
\]

%%%%%%%%%%%%%%%%%%%%%%%%%%%%%%%%%%%%%%%%%%%%%%%%%%%%%%%%%%%%
\subsection{Basic Laplace Transform Table}
\label{subsec:basic-laplace-table}
%%%%%%%%%%%%%%%%%%%%%%%%%%%%%%%%%%%%%%%%%%%%%%%%%%%%%%%%%%%%

\[
\boxed{
\begin{array}{c|c}
f(t) & \mathcal{L}\{f(t)\}
\\
\hline
1 & \dfrac1s
\\[0.8em]
t & \dfrac1{s^2}
\\[0.8em]
t^n & \dfrac{n!}{s^{n+1}}
\\[0.8em]
e^{at} & \dfrac1{s-a}
\\[0.8em]
\sin bt & \dfrac{b}{s^2+b^2}
\\[0.8em]
\cos bt & \dfrac{s}{s^2+b^2}
\\[0.8em]
\sinh at & \dfrac{a}{s^2-a^2}
\\[0.8em]
\cosh at & \dfrac{s}{s^2-a^2}
\end{array}
}
\]

%%%%%%%%%%%%%%%%%%%%%%%%%%%%%%%%%%%%%%%%%%%%%%%%%%%%%%%%%%%%
\subsection{Worked Example: Direct Laplace Transform}
\label{subsec:worked-direct-laplace}
%%%%%%%%%%%%%%%%%%%%%%%%%%%%%%%%%%%%%%%%%%%%%%%%%%%%%%%%%%%%

\begin{workedexamplebox}{Using Linearity}

Find

\[
\mathcal{L}
\{
3t^2-4\sin(2t)+5e^{-t}
\}.
\]

Use standard transforms:

\[
\mathcal{L}\{t^2\}
=
\frac2{s^3},
\]

\[
\mathcal{L}\{\sin 2t\}
=
\frac2{s^2+4},
\]

and

\[
\mathcal{L}\{e^{-t}\}
=
\frac1{s+1}.
\]

Therefore,

\begin{answerbox}
\[
\boxed{
\mathcal{L}
\{
3t^2-4\sin(2t)+5e^{-t}
\}
=
\frac6{s^3}
-
\frac8{s^2+4}
+
\frac5{s+1}.
}
\]
\end{answerbox}

\end{workedexamplebox}

%%%%%%%%%%%%%%%%%%%%%%%%%%%%%%%%%%%%%%%%%%%%%%%%%%%%%%%%%%%%
\subsection{First Shifting Theorem}
\label{subsec:first-shifting-theorem}
%%%%%%%%%%%%%%%%%%%%%%%%%%%%%%%%%%%%%%%%%%%%%%%%%%%%%%%%%%%%

If

\[
\mathcal{L}\{f(t)\}
=
F(s),
\]

then

\[
\boxed{
\mathcal{L}
\{
e^{at}f(t)
\}
=
F(s-a).
}
\]

This is sometimes called the exponential shift theorem.

%%%%%%%%%%%%%%%%%%%%%%%%%%%%%%%%%%%%%%%%%%%%%%%%%%%%%%%%%%%%
\subsection{Worked Example: First Shift}
\label{subsec:worked-first-shift}
%%%%%%%%%%%%%%%%%%%%%%%%%%%%%%%%%%%%%%%%%%%%%%%%%%%%%%%%%%%%

\begin{workedexamplebox}{Exponential Times Sine}

Find

\[
\mathcal{L}
\{
e^{3t}\sin(2t)
\}.
\]

Since

\[
\mathcal{L}\{\sin 2t\}
=
\frac2{s^2+4},
\]

replace \(s\) with

\[
s-3.
\]

Thus,

\begin{answerbox}
\[
\boxed{
\mathcal{L}
\{
e^{3t}\sin(2t)
\}
=
\frac2{
(s-3)^2+4
}.
}
\]
\end{answerbox}

\end{workedexamplebox}

%%%%%%%%%%%%%%%%%%%%%%%%%%%%%%%%%%%%%%%%%%%%%%%%%%%%%%%%%%%%
\subsection{Transforms of Derivatives}
\label{subsec:laplace-derivatives}
%%%%%%%%%%%%%%%%%%%%%%%%%%%%%%%%%%%%%%%%%%%%%%%%%%%%%%%%%%%%

Suppose

\[
F(s)
=
\mathcal{L}\{f(t)\}.
\]

Integration by parts gives

\[
\boxed{
\mathcal{L}\{f'(t)\}
=
sF(s)-f(0).
}
\]

This identity is one of the main reasons Laplace transforms are useful for
initial-value problems.

%%%%%%%%%%%%%%%%%%%%%%%%%%%%%%%%%%%%%%%%%%%%%%%%%%%%%%%%%%%%
\subsection{Second Derivative Transform}
\label{subsec:laplace-second-derivative}
%%%%%%%%%%%%%%%%%%%%%%%%%%%%%%%%%%%%%%%%%%%%%%%%%%%%%%%%%%%%

Applying the derivative formula again,

\[
\boxed{
\mathcal{L}\{f''(t)\}
=
s^2F(s)
-
sf(0)
-
f'(0).
}
\]

%%%%%%%%%%%%%%%%%%%%%%%%%%%%%%%%%%%%%%%%%%%%%%%%%%%%%%%%%%%%
\subsection{Higher Derivative Transform}
\label{subsec:laplace-higher-derivatives}
%%%%%%%%%%%%%%%%%%%%%%%%%%%%%%%%%%%%%%%%%%%%%%%%%%%%%%%%%%%%

For the \(n\)-th derivative,

\[
\boxed{
\mathcal{L}\{f^{(n)}(t)\}
=
s^nF(s)
-
s^{n-1}f(0)
-
s^{n-2}f'(0)
-
\cdots
-
f^{(n-1)}(0).
}
\]

Thus initial conditions appear automatically in the transformed equation.

%%%%%%%%%%%%%%%%%%%%%%%%%%%%%%%%%%%%%%%%%%%%%%%%%%%%%%%%%%%%
\subsection{Transform of an Integral}
\label{subsec:laplace-integral}
%%%%%%%%%%%%%%%%%%%%%%%%%%%%%%%%%%%%%%%%%%%%%%%%%%%%%%%%%%%%

If

\[
g(t)
=
\int_0^t
f(u)\,du,
\]

then

\[
g'(t)=f(t),
\qquad
g(0)=0.
\]

Therefore,

\[
sG(s)=F(s),
\]

so

\[
\boxed{
\mathcal{L}
\left\{
\int_0^t
f(u)\,du
\right\}
=
\frac{
F(s)
}{
s
}.
}
\]

%%%%%%%%%%%%%%%%%%%%%%%%%%%%%%%%%%%%%%%%%%%%%%%%%%%%%%%%%%%%
\subsection{Differentiation in the Transform Domain}
\label{subsec:differentiation-transform-domain}
%%%%%%%%%%%%%%%%%%%%%%%%%%%%%%%%%%%%%%%%%%%%%%%%%%%%%%%%%%%%

If

\[
F(s)
=
\mathcal{L}\{f(t)\},
\]

then differentiating with respect to \(s\) gives

\[
F'(s)
=
-
\int_0^\infty
t e^{-st}f(t)\,dt.
\]

Hence,

\[
\boxed{
\mathcal{L}\{tf(t)\}
=
-
F'(s).
}
\]

More generally,

\[
\boxed{
\mathcal{L}\{t^nf(t)\}
=
(-1)^n
F^{(n)}(s).
}
\]

%%%%%%%%%%%%%%%%%%%%%%%%%%%%%%%%%%%%%%%%%%%%%%%%%%%%%%%%%%%%
\subsection{Inverse Laplace Transform}
\label{subsec:inverse-laplace-transform}
%%%%%%%%%%%%%%%%%%%%%%%%%%%%%%%%%%%%%%%%%%%%%%%%%%%%%%%%%%%%

If

\[
F(s)
=
\mathcal{L}\{f(t)\},
\]

then

\[
\boxed{
f(t)
=
\mathcal{L}^{-1}\{F(s)\}.
}
\]

In elementary ODE problems, inverse transforms are usually found using:

\[
\boxed{
\text{transform tables},
}
\]

\[
\boxed{
\text{algebraic manipulation},
}
\]

\[
\boxed{
\text{partial fractions},
}
\]

and

\[
\boxed{
\text{shifting theorems}.
}
\]

%%%%%%%%%%%%%%%%%%%%%%%%%%%%%%%%%%%%%%%%%%%%%%%%%%%%%%%%%%%%
\subsection{Partial Fractions}
\label{subsec:laplace-partial-fractions}
%%%%%%%%%%%%%%%%%%%%%%%%%%%%%%%%%%%%%%%%%%%%%%%%%%%%%%%%%%%%

Suppose

\[
F(s)
=
\frac{
P(s)
}{
Q(s)
},
\]

where

\[
\deg P<\deg Q.
\]

Factoring \(Q(s)\) may allow decomposition into simpler rational terms.

For example,

\[
\boxed{
\frac{
3s+5
}{
(s+1)(s+2)
}
=
\frac{A}{s+1}
+
\frac{B}{s+2}.
}
\]

Each term can then be inverted separately.

%%%%%%%%%%%%%%%%%%%%%%%%%%%%%%%%%%%%%%%%%%%%%%%%%%%%%%%%%%%%
\subsection{Worked Example: Inverse Transform by Partial Fractions}
\label{subsec:worked-inverse-partial-fractions}
%%%%%%%%%%%%%%%%%%%%%%%%%%%%%%%%%%%%%%%%%%%%%%%%%%%%%%%%%%%%

\begin{workedexamplebox}{Distinct Linear Factors}

Find

\[
\mathcal{L}^{-1}
\left\{
\frac{
3s+5
}{
(s+1)(s+2)
}
\right\}.
\]

Write

\[
\frac{
3s+5
}{
(s+1)(s+2)
}
=
\frac{A}{s+1}
+
\frac{B}{s+2}.
\]

Then

\[
3s+5
=
A(s+2)
+
B(s+1).
\]

Matching coefficients gives

\[
A+B=3,
\]

\[
2A+B=5.
\]

Thus

\[
A=2,
\qquad
B=1.
\]

Therefore,

\begin{answerbox}
\[
\boxed{
f(t)
=
2e^{-t}
+
e^{-2t}.
}
\]
\end{answerbox}

\end{workedexamplebox}

%%%%%%%%%%%%%%%%%%%%%%%%%%%%%%%%%%%%%%%%%%%%%%%%%%%%%%%%%%%%
\subsection{Repeated Linear Factors}
\label{subsec:laplace-repeated-factors}
%%%%%%%%%%%%%%%%%%%%%%%%%%%%%%%%%%%%%%%%%%%%%%%%%%%%%%%%%%%%

If the denominator contains

\[
(s-a)^m,
\]

the partial fraction decomposition includes

\[
\boxed{
\frac{A_1}{s-a}
+
\frac{A_2}{(s-a)^2}
+
\cdots
+
\frac{A_m}{(s-a)^m}.
}
\]

These correspond to terms involving

\[
\boxed{
e^{at},
\quad
te^{at},
\quad
t^2e^{at},
\ldots
}
\]

in the inverse transform.

%%%%%%%%%%%%%%%%%%%%%%%%%%%%%%%%%%%%%%%%%%%%%%%%%%%%%%%%%%%%
\subsection{Irreducible Quadratic Factors}
\label{subsec:laplace-quadratic-factors}
%%%%%%%%%%%%%%%%%%%%%%%%%%%%%%%%%%%%%%%%%%%%%%%%%%%%%%%%%%%%

For a quadratic factor such as

\[
s^2+\omega^2,
\]

a partial fraction numerator has the form

\[
\boxed{
As+B.
}
\]

One then matches terms with

\[
\frac{s}{s^2+\omega^2}
\]

and

\[
\frac{\omega}{s^2+\omega^2}.
\]

These correspond to cosine and sine.

%%%%%%%%%%%%%%%%%%%%%%%%%%%%%%%%%%%%%%%%%%%%%%%%%%%%%%%%%%%%
\subsection{Completing the Square}
\label{subsec:laplace-completing-square}
%%%%%%%%%%%%%%%%%%%%%%%%%%%%%%%%%%%%%%%%%%%%%%%%%%%%%%%%%%%%

Expressions such as

\[
s^2-4s+13
\]

should be rewritten as

\[
\boxed{
(s-2)^2+9.
}
\]

Then first-shift formulas identify inverse transforms involving

\[
e^{2t}\cos3t
\]

or

\[
e^{2t}\sin3t.
\]

%%%%%%%%%%%%%%%%%%%%%%%%%%%%%%%%%%%%%%%%%%%%%%%%%%%%%%%%%%%%
\subsection{Solving Initial-Value Problems}
\label{subsec:laplace-solving-ivps}
%%%%%%%%%%%%%%%%%%%%%%%%%%%%%%%%%%%%%%%%%%%%%%%%%%%%%%%%%%%%

To solve an IVP using Laplace transforms:

\begin{enumerate}
    \item take the Laplace transform of both sides;
    \item insert the initial conditions;
    \item solve algebraically for \(Y(s)\);
    \item simplify \(Y(s)\);
    \item use partial fractions or transform identities;
    \item apply the inverse Laplace transform.
\end{enumerate}

%%%%%%%%%%%%%%%%%%%%%%%%%%%%%%%%%%%%%%%%%%%%%%%%%%%%%%%%%%%%
\subsection{Worked Example: First-Order IVP}
\label{subsec:worked-laplace-first-order-ivp}
%%%%%%%%%%%%%%%%%%%%%%%%%%%%%%%%%%%%%%%%%%%%%%%%%%%%%%%%%%%%

\begin{workedexamplebox}{First-Order Initial-Value Problem}

Solve

\[
y'+2y=4,
\qquad
y(0)=1.
\]

Let

\[
Y(s)
=
\mathcal{L}\{y(t)\}.
\]

Transform:

\[
sY(s)-1
+
2Y(s)
=
\frac4s.
\]

Thus,

\[
(s+2)Y(s)
=
1+\frac4s.
\]

Therefore,

\[
Y(s)
=
\frac{
s+4
}{
s(s+2)
}.
\]

Partial fractions give

\[
Y(s)
=
\frac2s
-
\frac1{s+2}.
\]

Hence,

\begin{answerbox}
\[
\boxed{
y(t)
=
2-e^{-2t}.
}
\]
\end{answerbox}

\end{workedexamplebox}

%%%%%%%%%%%%%%%%%%%%%%%%%%%%%%%%%%%%%%%%%%%%%%%%%%%%%%%%%%%%
\subsection{Worked Example: Second-Order IVP}
\label{subsec:worked-laplace-second-order-ivp}
%%%%%%%%%%%%%%%%%%%%%%%%%%%%%%%%%%%%%%%%%%%%%%%%%%%%%%%%%%%%

\begin{workedexamplebox}{Oscillator Initial-Value Problem}

Solve

\[
y''+4y=0,
\]

with

\[
y(0)=3,
\qquad
y'(0)=2.
\]

Transform:

\[
s^2Y(s)
-
3s
-
2
+
4Y(s)
=
0.
\]

Thus,

\[
Y(s)
=
\frac{
3s+2
}{
s^2+4
}.
\]

Separate:

\[
Y(s)
=
3
\frac{s}{s^2+4}
+
\frac2{s^2+4}.
\]

Since

\[
\mathcal{L}^{-1}
\left\{
\frac{s}{s^2+4}
\right\}
=
\cos2t,
\]

and

\[
\mathcal{L}^{-1}
\left\{
\frac2{s^2+4}
\right\}
=
\sin2t,
\]

we obtain

\begin{answerbox}
\[
\boxed{
y(t)
=
3\cos2t
+
\sin2t.
}
\]
\end{answerbox}

\end{workedexamplebox}

%%%%%%%%%%%%%%%%%%%%%%%%%%%%%%%%%%%%%%%%%%%%%%%%%%%%%%%%%%%%
\subsection{Piecewise Functions}
\label{subsec:laplace-piecewise-functions}
%%%%%%%%%%%%%%%%%%%%%%%%%%%%%%%%%%%%%%%%%%%%%%%%%%%%%%%%%%%%

One major advantage of Laplace transforms is that discontinuous forcing
functions can be represented algebraically using step functions.

This avoids solving the ODE separately on many intervals.

%%%%%%%%%%%%%%%%%%%%%%%%%%%%%%%%%%%%%%%%%%%%%%%%%%%%%%%%%%%%
\subsection{Heaviside Step Function}
\label{subsec:heaviside-step-function}
%%%%%%%%%%%%%%%%%%%%%%%%%%%%%%%%%%%%%%%%%%%%%%%%%%%%%%%%%%%%

\begin{definitionbox}{Heaviside Step Function}
The delayed unit step at \(t=a\) is

\[
\boxed{
u_a(t)
=
\begin{cases}
0,
&
t<a,
\\
1,
&
t\geq a.
\end{cases}
}
\]
\end{definitionbox}

The notation

\[
u(t-a)
\]

is also common.

%%%%%%%%%%%%%%%%%%%%%%%%%%%%%%%%%%%%%%%%%%%%%%%%%%%%%%%%%%%%
\subsection{Transform of the Unit Step}
\label{subsec:laplace-unit-step}
%%%%%%%%%%%%%%%%%%%%%%%%%%%%%%%%%%%%%%%%%%%%%%%%%%%%%%%%%%%%

For \(a\geq0\),

\[
\mathcal{L}\{u_a(t)\}
=
\int_a^\infty
e^{-st}\,dt.
\]

Therefore,

\[
\boxed{
\mathcal{L}\{u_a(t)\}
=
\frac{
e^{-as}
}{
s
}.
}
\]

%%%%%%%%%%%%%%%%%%%%%%%%%%%%%%%%%%%%%%%%%%%%%%%%%%%%%%%%%%%%
\subsection{Representing Piecewise Functions}
\label{subsec:step-function-representation}
%%%%%%%%%%%%%%%%%%%%%%%%%%%%%%%%%%%%%%%%%%%%%%%%%%%%%%%%%%%%

Suppose

\[
f(t)
=
\begin{cases}
f_1(t),
&
0\leq t<a,
\\
f_2(t),
&
t\geq a.
\end{cases}
\]

Then

\[
\boxed{
f(t)
=
f_1(t)
+
u_a(t)
[
f_2(t)-f_1(t)
].
}
\]

This representation makes the change at \(t=a\) explicit.

%%%%%%%%%%%%%%%%%%%%%%%%%%%%%%%%%%%%%%%%%%%%%%%%%%%%%%%%%%%%
\subsection{Worked Example: Step Representation}
\label{subsec:worked-step-representation}
%%%%%%%%%%%%%%%%%%%%%%%%%%%%%%%%%%%%%%%%%%%%%%%%%%%%%%%%%%%%

\begin{workedexamplebox}{Turning on a Constant}

Suppose

\[
f(t)
=
\begin{cases}
0,
&
0\leq t<3,
\\
5,
&
t\geq3.
\end{cases}
\]

Then

\[
\begin{answerbox}
\[
\boxed{
f(t)
=
5u_3(t).
}
\]
\end{answerbox}

Its Laplace transform is

\[
\boxed{
F(s)
=
\frac{
5e^{-3s}
}{
s
}.
}
\]

\end{workedexamplebox}

%%%%%%%%%%%%%%%%%%%%%%%%%%%%%%%%%%%%%%%%%%%%%%%%%%%%%%%%%%%%
\subsection{Second Shifting Theorem}
\label{subsec:second-shifting-theorem}
%%%%%%%%%%%%%%%%%%%%%%%%%%%%%%%%%%%%%%%%%%%%%%%%%%%%%%%%%%%%

If

\[
\mathcal{L}\{f(t)\}
=
F(s),
\]

then

\[
\boxed{
\mathcal{L}
\{
u_a(t)
f(t-a)
\}
=
e^{-as}
F(s).
}
\]

Equivalently,

\[
\boxed{
\mathcal{L}^{-1}
\{
e^{-as}F(s)
\}
=
u_a(t)
f(t-a).
}
\]

This is the time-delay theorem.

%%%%%%%%%%%%%%%%%%%%%%%%%%%%%%%%%%%%%%%%%%%%%%%%%%%%%%%%%%%%
\subsection{Worked Example: Delayed Function}
\label{subsec:worked-delayed-function}
%%%%%%%%%%%%%%%%%%%%%%%%%%%%%%%%%%%%%%%%%%%%%%%%%%%%%%%%%%%%

\begin{workedexamplebox}{Delayed Sine}

Find

\[
\mathcal{L}^{-1}
\left\{
e^{-4s}
\frac{
2
}{
s^2+4
}
\right\}.
\]

Since

\[
\mathcal{L}^{-1}
\left\{
\frac2{s^2+4}
\right\}
=
\sin2t,
\]

the second shifting theorem gives

\begin{answerbox}
\[
\boxed{
u_4(t)
\sin
\left[
2(t-4)
\right].
}
\]
\end{answerbox}

\end{workedexamplebox}

%%%%%%%%%%%%%%%%%%%%%%%%%%%%%%%%%%%%%%%%%%%%%%%%%%%%%%%%%%%%
\subsection{Delayed Forcing in ODEs}
\label{subsec:delayed-forcing-odes}
%%%%%%%%%%%%%%%%%%%%%%%%%%%%%%%%%%%%%%%%%%%%%%%%%%%%%%%%%%%%

Suppose an external force begins at time \(a\).

A forcing term can be written

\[
\boxed{
u_a(t)f(t-a).
}
\]

Its transform is immediately

\[
\boxed{
e^{-as}F(s).
}
\]

Thus delayed inputs become exponential factors in the transform domain.

%%%%%%%%%%%%%%%%%%%%%%%%%%%%%%%%%%%%%%%%%%%%%%%%%%%%%%%%%%%%
\subsection{Periodic Functions}
\label{subsec:laplace-periodic-functions}
%%%%%%%%%%%%%%%%%%%%%%%%%%%%%%%%%%%%%%%%%%%%%%%%%%%%%%%%%%%%

Suppose \(f(t)\) is periodic with period \(T>0\):

\[
f(t+T)=f(t).
\]

Then

\[
\boxed{
\mathcal{L}\{f(t)\}
=
\frac{
\displaystyle
\int_0^T
e^{-st}f(t)\,dt
}{
1-e^{-sT}
}.
}
\]

This is useful for repeated pulse, square-wave, and sawtooth forcing.

%%%%%%%%%%%%%%%%%%%%%%%%%%%%%%%%%%%%%%%%%%%%%%%%%%%%%%%%%%%%
\subsection{Dirac Delta Distribution}
\label{subsec:dirac-delta}
%%%%%%%%%%%%%%%%%%%%%%%%%%%%%%%%%%%%%%%%%%%%%%%%%%%%%%%%%%%%

The Dirac delta is not an ordinary function.

It is a generalized function or distribution used to model an idealized
instantaneous impulse.

It is written

\[
\boxed{
\delta(t-a).
}
\]

Conceptually, it is concentrated entirely at \(t=a\) while having unit
total mass.

%%%%%%%%%%%%%%%%%%%%%%%%%%%%%%%%%%%%%%%%%%%%%%%%%%%%%%%%%%%%
\subsection{Sifting Property}
\label{subsec:delta-sifting}
%%%%%%%%%%%%%%%%%%%%%%%%%%%%%%%%%%%%%%%%%%%%%%%%%%%%%%%%%%%%

The defining distributional identity is

\[
\boxed{
\int_{-\infty}^{\infty}
f(t)
\delta(t-a)
\,dt
=
f(a),
}
\]

for suitable test functions \(f\).

On the half-line with \(a>0\),

\[
\boxed{
\int_0^\infty
f(t)
\delta(t-a)
\,dt
=
f(a).
}
\]

%%%%%%%%%%%%%%%%%%%%%%%%%%%%%%%%%%%%%%%%%%%%%%%%%%%%%%%%%%%%
\subsection{Laplace Transform of a Delta Impulse}
\label{subsec:laplace-delta}
%%%%%%%%%%%%%%%%%%%%%%%%%%%%%%%%%%%%%%%%%%%%%%%%%%%%%%%%%%%%

Using the sifting property,

\[
\mathcal{L}
\{
\delta(t-a)
\}
=
\int_0^\infty
e^{-st}
\delta(t-a)
\,dt.
\]

Therefore,

\[
\boxed{
\mathcal{L}
\{
\delta(t-a)
\}
=
e^{-as},
\qquad
a>0.
}
\]

For an impulse at the origin under the standard one-sided transform
convention,

\[
\boxed{
\mathcal{L}\{\delta(t)\}=1.
}
\]

%%%%%%%%%%%%%%%%%%%%%%%%%%%%%%%%%%%%%%%%%%%%%%%%%%%%%%%%%%%%
\subsection{Impulse Forcing}
\label{subsec:impulse-forcing}
%%%%%%%%%%%%%%%%%%%%%%%%%%%%%%%%%%%%%%%%%%%%%%%%%%%%%%%%%%%%

An idealized impulsive force may be written

\[
\boxed{
F_0
\delta(t-a).
}
\]

Its Laplace transform is

\[
\boxed{
F_0e^{-as}.
}
\]

This allows differential equations with sudden impacts to be solved
algebraically.

%%%%%%%%%%%%%%%%%%%%%%%%%%%%%%%%%%%%%%%%%%%%%%%%%%%%%%%%%%%%
\subsection{Worked Example: Impulse Response}
\label{subsec:worked-impulse-response}
%%%%%%%%%%%%%%%%%%%%%%%%%%%%%%%%%%%%%%%%%%%%%%%%%%%%%%%%%%%%

\begin{workedexamplebox}{Impulse Applied to an Oscillator}

Consider

\[
y''+y
=
\delta(t-\pi),
\]

with

\[
y(0)=0,
\qquad
y'(0)=0.
\]

Transform:

\[
(s^2+1)Y(s)
=
e^{-\pi s}.
\]

Thus

\[
Y(s)
=
e^{-\pi s}
\frac1{s^2+1}.
\]

Since

\[
\mathcal{L}^{-1}
\left\{
\frac1{s^2+1}
\right\}
=
\sin t,
\]

we obtain

\begin{answerbox}
\[
\boxed{
y(t)
=
u_\pi(t)
\sin(t-\pi).
}
\]
\end{answerbox}

\end{workedexamplebox}

%%%%%%%%%%%%%%%%%%%%%%%%%%%%%%%%%%%%%%%%%%%%%%%%%%%%%%%%%%%%
\subsection{Convolution}
\label{subsec:convolution}
%%%%%%%%%%%%%%%%%%%%%%%%%%%%%%%%%%%%%%%%%%%%%%%%%%%%%%%%%%%%

\begin{definitionbox}{Convolution}
For functions \(f\) and \(g\) defined on \(t\geq0\), their convolution is

\[
\boxed{
(f*g)(t)
=
\int_0^t
f(\tau)
g(t-\tau)
\,d\tau.
}
\]
\end{definitionbox}

Convolution combines two functions through shifted overlap.

%%%%%%%%%%%%%%%%%%%%%%%%%%%%%%%%%%%%%%%%%%%%%%%%%%%%%%%%%%%%
\subsection{Commutativity of Convolution}
\label{subsec:convolution-commutative}
%%%%%%%%%%%%%%%%%%%%%%%%%%%%%%%%%%%%%%%%%%%%%%%%%%%%%%%%%%%%

Convolution satisfies

\[
\boxed{
f*g
=
g*f.
}
\]

Indeed,

\[
\int_0^t
f(\tau)
g(t-\tau)
\,d\tau
\]

can be transformed by the substitution

\[
u=t-\tau.
\]

%%%%%%%%%%%%%%%%%%%%%%%%%%%%%%%%%%%%%%%%%%%%%%%%%%%%%%%%%%%%
\subsection{Other Convolution Properties}
\label{subsec:convolution-properties}
%%%%%%%%%%%%%%%%%%%%%%%%%%%%%%%%%%%%%%%%%%%%%%%%%%%%%%%%%%%%

Under suitable conditions,

\[
\boxed{
f*(g+h)
=
f*g+f*h,
}
\]

\[
\boxed{
(f*g)*h
=
f*(g*h),
}
\]

and

\[
\boxed{
cf*g
=
c(f*g).
}
\]

Thus convolution is bilinear, associative, and commutative.

%%%%%%%%%%%%%%%%%%%%%%%%%%%%%%%%%%%%%%%%%%%%%%%%%%%%%%%%%%%%
\subsection{Convolution Theorem}
\label{subsec:laplace-convolution-theorem}
%%%%%%%%%%%%%%%%%%%%%%%%%%%%%%%%%%%%%%%%%%%%%%%%%%%%%%%%%%%%

If

\[
\mathcal{L}\{f\}=F
\]

and

\[
\mathcal{L}\{g\}=G,
\]

then

\[
\boxed{
\mathcal{L}\{f*g\}
=
F(s)G(s).
}
\]

Therefore,

\[
\boxed{
\mathcal{L}^{-1}
\{
F(s)G(s)
\}
=
(f*g)(t).
}
\]

%%%%%%%%%%%%%%%%%%%%%%%%%%%%%%%%%%%%%%%%%%%%%%%%%%%%%%%%%%%%
\subsection{Worked Example: Convolution}
\label{subsec:worked-convolution}
%%%%%%%%%%%%%%%%%%%%%%%%%%%%%%%%%%%%%%%%%%%%%%%%%%%%%%%%%%%%

\begin{workedexamplebox}{Inverse Transform Using Convolution}

Find

\[
\mathcal{L}^{-1}
\left\{
\frac1{
s(s^2+1)
}
\right\}.
\]

Since

\[
\frac1s
=
\mathcal{L}\{1\},
\]

and

\[
\frac1{s^2+1}
=
\mathcal{L}\{\sin t\},
\]

the inverse transform is

\[
(1*\sin)(t)
=
\int_0^t
\sin(t-\tau)\,d\tau.
\]

Thus,

\[
\int_0^t
\sin(t-\tau)\,d\tau
=
1-\cos t.
\]

Therefore,

\begin{answerbox}
\[
\boxed{
\mathcal{L}^{-1}
\left\{
\frac1{s(s^2+1)}
\right\}
=
1-\cos t.
}
\]
\end{answerbox}

\end{workedexamplebox}

%%%%%%%%%%%%%%%%%%%%%%%%%%%%%%%%%%%%%%%%%%%%%%%%%%%%%%%%%%%%
\subsection{Impulse Response}
\label{subsec:impulse-response}
%%%%%%%%%%%%%%%%%%%%%%%%%%%%%%%%%%%%%%%%%%%%%%%%%%%%%%%%%%%%

Consider a linear time-invariant differential equation with zero initial
conditions.

If an impulse

\[
\delta(t)
\]

is applied as input, the resulting output is called the impulse response.

Denote it by

\[
\boxed{
h(t).
}
\]

Its Laplace transform is the system transfer function under standard
zero-initial-condition conventions.

%%%%%%%%%%%%%%%%%%%%%%%%%%%%%%%%%%%%%%%%%%%%%%%%%%%%%%%%%%%%
\subsection{Transfer Function}
\label{subsec:transfer-function}
%%%%%%%%%%%%%%%%%%%%%%%%%%%%%%%%%%%%%%%%%%%%%%%%%%%%%%%%%%%%

Suppose a linear system relates input \(u(t)\) and output \(y(t)\).

Under zero initial conditions, define

\[
\boxed{
H(s)
=
\frac{
Y(s)
}{
U(s)
}.
}
\]

The function \(H(s)\) is the transfer function.

%%%%%%%%%%%%%%%%%%%%%%%%%%%%%%%%%%%%%%%%%%%%%%%%%%%%%%%%%%%%
\subsection{Example of a Transfer Function}
\label{subsec:transfer-function-example}
%%%%%%%%%%%%%%%%%%%%%%%%%%%%%%%%%%%%%%%%%%%%%%%%%%%%%%%%%%%%

Consider

\[
\boxed{
y''+3y'+2y=u(t).
}
\]

Under zero initial conditions,

\[
(s^2+3s+2)Y(s)
=
U(s).
\]

Therefore,

\[
\boxed{
H(s)
=
\frac{
Y(s)
}{
U(s)
}
=
\frac1{
s^2+3s+2
}.
}
\]

Since

\[
s^2+3s+2
=
(s+1)(s+2),
\]

the poles occur at

\[
\boxed{
s=-1,
\qquad
s=-2.
}
\]

%%%%%%%%%%%%%%%%%%%%%%%%%%%%%%%%%%%%%%%%%%%%%%%%%%%%%%%%%%%%
\subsection{Output as Convolution}
\label{subsec:output-convolution}
%%%%%%%%%%%%%%%%%%%%%%%%%%%%%%%%%%%%%%%%%%%%%%%%%%%%%%%%%%%%

If

\[
Y(s)
=
H(s)U(s),
\]

then the convolution theorem gives

\[
\boxed{
y(t)
=
(h*u)(t).
}
\]

Therefore,

\[
\boxed{
y(t)
=
\int_0^t
h(\tau)
u(t-\tau)
\,d\tau.
}
\]

The impulse response completely characterizes a linear time-invariant
system under suitable conditions.

%%%%%%%%%%%%%%%%%%%%%%%%%%%%%%%%%%%%%%%%%%%%%%%%%%%%%%%%%%%%
\subsection{Step Response}
\label{subsec:step-response}
%%%%%%%%%%%%%%%%%%%%%%%%%%%%%%%%%%%%%%%%%%%%%%%%%%%%%%%%%%%%

If the input is the unit step

\[
u(t)=1,
\]

then

\[
U(s)
=
\frac1s.
\]

The output is

\[
\boxed{
Y(s)
=
\frac{
H(s)
}{
s
}.
}
\]

The corresponding inverse transform is called the step response.

%%%%%%%%%%%%%%%%%%%%%%%%%%%%%%%%%%%%%%%%%%%%%%%%%%%%%%%%%%%%
\subsection{Poles}
\label{subsec:transform-poles}
%%%%%%%%%%%%%%%%%%%%%%%%%%%%%%%%%%%%%%%%%%%%%%%%%%%%%%%%%%%%

A pole is a value of \(s\) where a transform becomes unbounded due to a
zero of its denominator after cancellation.

For a rational transfer function

\[
\boxed{
H(s)
=
\frac{
N(s)
}{
D(s)
},
}
\]

poles are roots of

\[
\boxed{
D(s)=0.
}
\]

%%%%%%%%%%%%%%%%%%%%%%%%%%%%%%%%%%%%%%%%%%%%%%%%%%%%%%%%%%%%
\subsection{Poles and Differential Equations}
\label{subsec:poles-differential-equations}
%%%%%%%%%%%%%%%%%%%%%%%%%%%%%%%%%%%%%%%%%%%%%%%%%%%%%%%%%%%%

For constant-coefficient linear ODEs, the denominator polynomial in the
Laplace domain is closely related to the characteristic polynomial.

Thus characteristic roots in ODE theory correspond to poles in transform
theory.

For example,

\[
y''+3y'+2y=0
\]

has characteristic roots

\[
-1,
\qquad
-2,
\]

matching poles of

\[
\frac1{(s+1)(s+2)}.
\]

%%%%%%%%%%%%%%%%%%%%%%%%%%%%%%%%%%%%%%%%%%%%%%%%%%%%%%%%%%%%
\subsection{Pole Location and Stability}
\label{subsec:pole-location-stability}
%%%%%%%%%%%%%%%%%%%%%%%%%%%%%%%%%%%%%%%%%%%%%%%%%%%%%%%%%%%%

For a finite-dimensional continuous-time linear system, poles with

\[
\boxed{
\operatorname{Re}(s)<0
}
\]

correspond to exponentially decaying modes.

Poles with

\[
\boxed{
\operatorname{Re}(s)>0
}
\]

correspond to exponentially growing modes.

Poles on the imaginary axis require more careful analysis, especially
when repeated.

%%%%%%%%%%%%%%%%%%%%%%%%%%%%%%%%%%%%%%%%%%%%%%%%%%%%%%%%%%%%
\subsection{Zeros}
\label{subsec:transform-zeros}
%%%%%%%%%%%%%%%%%%%%%%%%%%%%%%%%%%%%%%%%%%%%%%%%%%%%%%%%%%%%

For

\[
H(s)
=
\frac{
N(s)
}{
D(s)
},
\]

zeros are roots of

\[
\boxed{
N(s)=0.
}
\]

Zeros influence how particular input frequencies or modes appear in the
output.

Pole-zero structure is central to control and signal analysis.

%%%%%%%%%%%%%%%%%%%%%%%%%%%%%%%%%%%%%%%%%%%%%%%%%%%%%%%%%%%%
\subsection{Solving Systems with Laplace Transforms}
\label{subsec:laplace-systems}
%%%%%%%%%%%%%%%%%%%%%%%%%%%%%%%%%%%%%%%%%%%%%%%%%%%%%%%%%%%%

Consider

\[
\boxed{
\mathbf{x}'
=
A\mathbf{x}
+
\mathbf{f}(t),
}
\]

with

\[
\mathbf{x}(0)
=
\mathbf{x}_0.
\]

Taking transforms gives

\[
s\mathbf{X}(s)
-
\mathbf{x}_0
=
A\mathbf{X}(s)
+
\mathbf{F}(s).
\]

Thus,

\[
\boxed{
(
sI-A
)
\mathbf{X}(s)
=
\mathbf{x}_0
+
\mathbf{F}(s).
}
\]

%%%%%%%%%%%%%%%%%%%%%%%%%%%%%%%%%%%%%%%%%%%%%%%%%%%%%%%%%%%%
\subsection{Resolvent Matrix}
\label{subsec:resolvent-matrix}
%%%%%%%%%%%%%%%%%%%%%%%%%%%%%%%%%%%%%%%%%%%%%%%%%%%%%%%%%%%%

If

\[
sI-A
\]

is invertible, then

\[
\boxed{
\mathbf{X}(s)
=
(
sI-A
)^{-1}
[
\mathbf{x}_0+\mathbf{F}(s)
].
}
\]

The matrix

\[
\boxed{
(
sI-A
)^{-1}
}
\]

is called the resolvent of \(A\).

It connects transform methods with the matrix exponential.

%%%%%%%%%%%%%%%%%%%%%%%%%%%%%%%%%%%%%%%%%%%%%%%%%%%%%%%%%%%%
\subsection{Matrix Exponential and Laplace Transform}
\label{subsec:matrix-exponential-laplace}
%%%%%%%%%%%%%%%%%%%%%%%%%%%%%%%%%%%%%%%%%%%%%%%%%%%%%%%%%%%%

For suitable \(s\),

\[
\boxed{
\mathcal{L}
\{
e^{At}
\}
=
(
sI-A
)^{-1}.
}
\]

Thus the matrix exponential is the inverse Laplace transform of the
resolvent:

\[
\boxed{
e^{At}
=
\mathcal{L}^{-1}
\left\{
(
sI-A
)^{-1}
\right\}.
}
\]

%%%%%%%%%%%%%%%%%%%%%%%%%%%%%%%%%%%%%%%%%%%%%%%%%%%%%%%%%%%%
\subsection{Variation of Constants from Convolution}
\label{subsec:variation-constants-transform}
%%%%%%%%%%%%%%%%%%%%%%%%%%%%%%%%%%%%%%%%%%%%%%%%%%%%%%%%%%%%

For

\[
\mathbf{x}'
=
A\mathbf{x}
+
\mathbf{f}(t),
\]

the transformed solution implies

\[
\boxed{
\mathbf{x}(t)
=
e^{At}\mathbf{x}_0
+
\int_0^t
e^{A(t-\tau)}
\mathbf{f}(\tau)
\,d\tau.
}
\]

This is the variation-of-constants formula.

The integral term is a convolution with the matrix impulse response.

%%%%%%%%%%%%%%%%%%%%%%%%%%%%%%%%%%%%%%%%%%%%%%%%%%%%%%%%%%%%
\subsection{Initial Value Theorem}
\label{subsec:initial-value-theorem}
%%%%%%%%%%%%%%%%%%%%%%%%%%%%%%%%%%%%%%%%%%%%%%%%%%%%%%%%%%%%

Under suitable regularity conditions,

\[
\boxed{
f(0^+)
=
\lim_{s\to\infty}
sF(s).
}
\]

This allows the initial value to be recovered directly from the transform.

It should be applied only when its hypotheses are satisfied.

%%%%%%%%%%%%%%%%%%%%%%%%%%%%%%%%%%%%%%%%%%%%%%%%%%%%%%%%%%%%
\subsection{Final Value Theorem}
\label{subsec:final-value-theorem}
%%%%%%%%%%%%%%%%%%%%%%%%%%%%%%%%%%%%%%%%%%%%%%%%%%%%%%%%%%%%

Under appropriate stability conditions,

\[
\boxed{
\lim_{t\to\infty}
f(t)
=
\lim_{s\to0}
sF(s).
}
\]

The theorem is not valid for arbitrary functions.

For rational systems, poles of \(sF(s)\) must satisfy appropriate
left-half-plane conditions, with the possible pole at \(s=0\) handled
carefully.

%%%%%%%%%%%%%%%%%%%%%%%%%%%%%%%%%%%%%%%%%%%%%%%%%%%%%%%%%%%%
\subsection{Worked Example: Final Value}
\label{subsec:worked-final-value}
%%%%%%%%%%%%%%%%%%%%%%%%%%%%%%%%%%%%%%%%%%%%%%%%%%%%%%%%%%%%

\begin{workedexamplebox}{Steady-State Value}

Suppose

\[
F(s)
=
\frac{
2
}{
s(s+4)
}.
\]

Then

\[
sF(s)
=
\frac2{s+4}.
\]

The pole of \(sF(s)\) is

\[
s=-4,
\]

so the standard stability condition is satisfied.

Therefore,

\begin{answerbox}
\[
\boxed{
\lim_{t\to\infty}
f(t)
=
\lim_{s\to0}
\frac2{s+4}
=
\frac12.
}
\]
\end{answerbox}

\end{workedexamplebox}

%%%%%%%%%%%%%%%%%%%%%%%%%%%%%%%%%%%%%%%%%%%%%%%%%%%%%%%%%%%%
\subsection{Bilateral Laplace Transform Preview}
\label{subsec:bilateral-laplace-preview}
%%%%%%%%%%%%%%%%%%%%%%%%%%%%%%%%%%%%%%%%%%%%%%%%%%%%%%%%%%%%

The one-sided Laplace transform integrates over

\[
t\geq0.
\]

The bilateral Laplace transform uses the entire real line:

\[
\boxed{
F(s)
=
\int_{-\infty}^{\infty}
e^{-st}
f(t)\,dt.
}
\]

Its convergence generally occurs in a region of the complex \(s\)-plane
called the region of convergence.

%%%%%%%%%%%%%%%%%%%%%%%%%%%%%%%%%%%%%%%%%%%%%%%%%%%%%%%%%%%%
\subsection{Complex Transform Variable}
\label{subsec:complex-laplace-variable}
%%%%%%%%%%%%%%%%%%%%%%%%%%%%%%%%%%%%%%%%%%%%%%%%%%%%%%%%%%%%

More generally,

\[
\boxed{
s
=
\sigma+i\omega.
}
\]

Then

\[
e^{-st}
=
e^{-\sigma t}
e^{-i\omega t}.
\]

Thus the Laplace transform combines:

\[
\boxed{
\text{exponential weighting}
}
\]

with

\[
\boxed{
\text{frequency analysis}.
}
\]

This reveals its close connection to the Fourier transform.

%%%%%%%%%%%%%%%%%%%%%%%%%%%%%%%%%%%%%%%%%%%%%%%%%%%%%%%%%%%%
\subsection{Fourier Transform}
\label{subsec:fourier-transform}
%%%%%%%%%%%%%%%%%%%%%%%%%%%%%%%%%%%%%%%%%%%%%%%%%%%%%%%%%%%%

Using the convention adopted here, the Fourier transform of \(f\) is

\[
\boxed{
\widehat{f}(\omega)
=
\int_{-\infty}^{\infty}
f(t)e^{-i\omega t}\,dt.
}
\]

The inverse transform is

\[
\boxed{
f(t)
=
\frac1{2\pi}
\int_{-\infty}^{\infty}
\widehat{f}(\omega)
e^{i\omega t}
\,d\omega,
}
\]

under suitable hypotheses.

%%%%%%%%%%%%%%%%%%%%%%%%%%%%%%%%%%%%%%%%%%%%%%%%%%%%%%%%%%%%
\subsection{Laplace and Fourier Connection}
\label{subsec:laplace-fourier-connection}
%%%%%%%%%%%%%%%%%%%%%%%%%%%%%%%%%%%%%%%%%%%%%%%%%%%%%%%%%%%%

For

\[
s=\sigma+i\omega,
\]

the bilateral Laplace transform becomes

\[
\int_{-\infty}^{\infty}
f(t)
e^{-\sigma t}
e^{-i\omega t}
\,dt.
\]

Thus it is the Fourier transform of the exponentially weighted function

\[
\boxed{
e^{-\sigma t}f(t).
}
\]

The Fourier transform corresponds formally to evaluating the bilateral
Laplace transform on

\[
\boxed{
s=i\omega,
}
\]

when the relevant convergence conditions hold.

%%%%%%%%%%%%%%%%%%%%%%%%%%%%%%%%%%%%%%%%%%%%%%%%%%%%%%%%%%%%
\subsection{Fourier Transform of a Derivative}
\label{subsec:fourier-transform-derivative}
%%%%%%%%%%%%%%%%%%%%%%%%%%%%%%%%%%%%%%%%%%%%%%%%%%%%%%%%%%%%

Under sufficient decay and smoothness,

\[
\boxed{
\mathcal{F}\{f'(t)\}
=
i\omega
\widehat{f}(\omega).
}
\]

Similarly,

\[
\boxed{
\mathcal{F}\{f^{(n)}(t)\}
=
(i\omega)^n
\widehat{f}(\omega).
}
\]

Thus differentiation becomes multiplication by powers of \(i\omega\).

%%%%%%%%%%%%%%%%%%%%%%%%%%%%%%%%%%%%%%%%%%%%%%%%%%%%%%%%%%%%
\subsection{Fourier Convolution Theorem}
\label{subsec:fourier-convolution}
%%%%%%%%%%%%%%%%%%%%%%%%%%%%%%%%%%%%%%%%%%%%%%%%%%%%%%%%%%%%

With the Fourier convention used here,

\[
\boxed{
\mathcal{F}\{f*g\}
=
\widehat{f}(\omega)
\widehat{g}(\omega).
}
\]

Thus convolution in the original domain becomes multiplication in the
frequency domain.

This is fundamental in signal processing and PDE analysis.

%%%%%%%%%%%%%%%%%%%%%%%%%%%%%%%%%%%%%%%%%%%%%%%%%%%%%%%%%%%%
\subsection{Frequency Interpretation}
\label{subsec:frequency-interpretation}
%%%%%%%%%%%%%%%%%%%%%%%%%%%%%%%%%%%%%%%%%%%%%%%%%%%%%%%%%%%%

The Fourier transform decomposes a signal into frequency components.

Large values of

\[
|\widehat{f}(\omega)|
\]

indicate substantial contribution from frequency \(\omega\).

This allows differential equations and physical systems to be studied
through their frequency response.

%%%%%%%%%%%%%%%%%%%%%%%%%%%%%%%%%%%%%%%%%%%%%%%%%%%%%%%%%%%%
\subsection{Fourier Sine Transform}
\label{subsec:fourier-sine-transform}
%%%%%%%%%%%%%%%%%%%%%%%%%%%%%%%%%%%%%%%%%%%%%%%%%%%%%%%%%%%%

For functions on the positive half-line, a Fourier sine transform may be
defined by

\[
\boxed{
F_s(\omega)
=
\int_0^\infty
f(x)
\sin(\omega x)
\,dx.
}
\]

Normalization conventions vary.

Sine transforms naturally appear with certain zero-value boundary
conditions.

%%%%%%%%%%%%%%%%%%%%%%%%%%%%%%%%%%%%%%%%%%%%%%%%%%%%%%%%%%%%
\subsection{Fourier Cosine Transform}
\label{subsec:fourier-cosine-transform}
%%%%%%%%%%%%%%%%%%%%%%%%%%%%%%%%%%%%%%%%%%%%%%%%%%%%%%%%%%%%

A Fourier cosine transform may be defined by

\[
\boxed{
F_c(\omega)
=
\int_0^\infty
f(x)
\cos(\omega x)
\,dx.
}
\]

Cosine transforms naturally arise with certain derivative-type boundary
conditions.

These transforms become especially important for PDEs.

%%%%%%%%%%%%%%%%%%%%%%%%%%%%%%%%%%%%%%%%%%%%%%%%%%%%%%%%%%%%
\subsection{Transform Methods for Differential Equations}
\label{subsec:transform-differential-equations}
%%%%%%%%%%%%%%%%%%%%%%%%%%%%%%%%%%%%%%%%%%%%%%%%%%%%%%%%%%%%

The general transform strategy is:

\[
\boxed{
\text{differentiate in original domain}
}
\]

becomes

\[
\boxed{
\text{multiply in transform domain}.
}
\]

For example,

\[
y''+ay'+by=f(t)
\]

becomes an algebraic equation involving

\[
Y(s).
\]

This eliminates direct differentiation of the unknown function.

%%%%%%%%%%%%%%%%%%%%%%%%%%%%%%%%%%%%%%%%%%%%%%%%%%%%%%%%%%%%
\subsection{Transform Methods and Boundary Problems}
\label{subsec:transform-boundary-problems}
%%%%%%%%%%%%%%%%%%%%%%%%%%%%%%%%%%%%%%%%%%%%%%%%%%%%%%%%%%%%

Fourier transforms and related methods are useful when differential
equations are posed on unbounded spatial domains.

Fourier sine and cosine transforms are useful on half-lines.

Laplace transforms are particularly natural for initial-value problems in
time.

Thus the transform is selected according to:

\[
\boxed{
\text{domain},
}
\]

\[
\boxed{
\text{boundary conditions},
}
\]

and

\[
\boxed{
\text{type of forcing}.
}
\]

%%%%%%%%%%%%%%%%%%%%%%%%%%%%%%%%%%%%%%%%%%%%%%%%%%%%%%%%%%%%
\subsection{The \(z\)-Transform Preview}
\label{subsec:z-transform-preview}
%%%%%%%%%%%%%%%%%%%%%%%%%%%%%%%%%%%%%%%%%%%%%%%%%%%%%%%%%%%%

For a discrete sequence

\[
x_0,x_1,x_2,\ldots,
\]

a one-sided \(z\)-transform is

\[
\boxed{
X(z)
=
\sum_{n=0}^{\infty}
x_nz^{-n}.
}
\]

The \(z\)-transform plays a role for difference equations analogous to
the Laplace transform for differential equations.

%%%%%%%%%%%%%%%%%%%%%%%%%%%%%%%%%%%%%%%%%%%%%%%%%%%%%%%%%%%%
\subsection{Shift Property of the \(z\)-Transform}
\label{subsec:z-transform-shift-preview}
%%%%%%%%%%%%%%%%%%%%%%%%%%%%%%%%%%%%%%%%%%%%%%%%%%%%%%%%%%%%

For a one-sided transform, shifting a sequence introduces algebraic
factors in \(z\) together with initial terms.

This allows recurrence relations to be converted into algebraic
equations.

The method is developed further when difference equations are studied
later in Chapter 9.

%%%%%%%%%%%%%%%%%%%%%%%%%%%%%%%%%%%%%%%%%%%%%%%%%%%%%%%%%%%%
\subsection{Laplace Versus Fourier Versus \(z\)-Transform}
\label{subsec:transform-comparison}
%%%%%%%%%%%%%%%%%%%%%%%%%%%%%%%%%%%%%%%%%%%%%%%%%%%%%%%%%%%%

A useful conceptual comparison is

\[
\begin{array}{c|c|c}
\text{Transform}
&
\text{Typical Domain}
&
\text{Common Use}
\\
\hline
\text{Laplace}
&
t\geq0
&
\text{continuous-time IVPs}
\\
\text{Fourier}
&
-\infty<t<\infty
&
\text{frequency analysis and PDEs}
\\
z\text{-transform}
&
n=0,1,2,\ldots
&
\text{difference equations}
\end{array}
\]

The boundaries between these applications are not absolute.

%%%%%%%%%%%%%%%%%%%%%%%%%%%%%%%%%%%%%%%%%%%%%%%%%%%%%%%%%%%%
\subsection{Transform Method Workflow}
\label{subsec:transform-method-workflow}
%%%%%%%%%%%%%%%%%%%%%%%%%%%%%%%%%%%%%%%%%%%%%%%%%%%%%%%%%%%%

For an ODE initial-value problem:

\begin{enumerate}
    \item identify the unknown function and initial conditions;
    \item transform both sides of the equation;
    \item replace derivative transforms using the initial data;
    \item solve algebraically for the transformed unknown;
    \item factor and simplify the resulting expression;
    \item use partial fractions if necessary;
    \item identify exponential shifts or time delays;
    \item apply convolution when products of transforms are convenient;
    \item take the inverse transform;
    \item verify the initial conditions;
    \item verify the differential equation on intervals where the
          ordinary equation applies;
    \item interpret discontinuities or impulses carefully.
\end{enumerate}

%%%%%%%%%%%%%%%%%%%%%%%%%%%%%%%%%%%%%%%%%%%%%%%%%%%%%%%%%%%%
\subsection{Choosing Transform Techniques}
\label{subsec:choosing-transform-technique}
%%%%%%%%%%%%%%%%%%%%%%%%%%%%%%%%%%%%%%%%%%%%%%%%%%%%%%%%%%%%

Use direct transform tables when:

\[
\boxed{
F(s)
\text{ matches a standard form}.
}
\]

Use partial fractions when:

\[
\boxed{
F(s)
\text{ is rational}.
}
\]

Use the first shifting theorem when:

\[
\boxed{
s
\text{ appears as }
s-a.
}
\]

Use the second shifting theorem when:

\[
\boxed{
e^{-as}
\text{ appears}.
}
\]

Use convolution when:

\[
\boxed{
F(s)
=
F_1(s)F_2(s)
}
\]

and direct inversion is inconvenient.

%%%%%%%%%%%%%%%%%%%%%%%%%%%%%%%%%%%%%%%%%%%%%%%%%%%%%%%%%%%%
\subsection{Common Errors}
\label{subsec:transform-method-common-errors}
%%%%%%%%%%%%%%%%%%%%%%%%%%%%%%%%%%%%%%%%%%%%%%%%%%%%%%%%%%%%

\subsubsection{Forgetting Initial Terms in Derivative Transforms}

The transform is not simply

\[
\mathcal{L}\{y'\}
=
sY.
\]

The correct formula is

\[
\boxed{
\mathcal{L}\{y'\}
=
sY-y(0).
}
\]

Similarly,

\[
\boxed{
\mathcal{L}\{y''\}
=
s^2Y
-
sy(0)
-
y'(0).
}
\]

%%%%%%%%%%%%%%%%%%%%%%%%%%%%%%%%%%%%%%%%%%%%%%%%%%%%%%%%%%%%
\subsubsection{Confusing the Two Shifting Theorems}

The first shift concerns multiplication by

\[
e^{at}.
\]

The second shift concerns a time delay represented by

\[
e^{-as}
\]

in the transform domain.

They solve different problems.

%%%%%%%%%%%%%%%%%%%%%%%%%%%%%%%%%%%%%%%%%%%%%%%%%%%%%%%%%%%%
\subsubsection{Forgetting the Shift Inside the Function}

The inverse transform

\[
e^{-as}F(s)
\]

is

\[
\boxed{
u_a(t)f(t-a),
}
\]

not simply

\[
u_a(t)f(t).
\]

%%%%%%%%%%%%%%%%%%%%%%%%%%%%%%%%%%%%%%%%%%%%%%%%%%%%%%%%%%%%
\subsubsection{Using a Step Function Without Rewriting Around the Delay}

For the second shifting theorem, expressions should usually be rewritten
in terms of

\[
\boxed{
t-a.
}
\]

This makes the transform rule directly applicable.

%%%%%%%%%%%%%%%%%%%%%%%%%%%%%%%%%%%%%%%%%%%%%%%%%%%%%%%%%%%%
\subsubsection{Treating the Dirac Delta as an Ordinary Function}

The Dirac delta is a distribution.

Statements such as

\[
\delta(0)=\infty
\]

should not be used as literal ordinary-function definitions.

Its meaning is given through integration against test functions.

%%%%%%%%%%%%%%%%%%%%%%%%%%%%%%%%%%%%%%%%%%%%%%%%%%%%%%%%%%%%
\subsubsection{Dropping Initial Conditions Before Transforming}

Initial conditions are a major advantage of Laplace methods.

They should be inserted directly into the transformed derivative
expressions.

%%%%%%%%%%%%%%%%%%%%%%%%%%%%%%%%%%%%%%%%%%%%%%%%%%%%%%%%%%%%
\subsubsection{Forgetting to Make a Rational Expression Proper}

If

\[
\deg P
\geq
\deg Q,
\]

perform polynomial division before ordinary partial-fraction
decomposition.

%%%%%%%%%%%%%%%%%%%%%%%%%%%%%%%%%%%%%%%%%%%%%%%%%%%%%%%%%%%%
\subsubsection{Incorrect Partial-Fraction Terms}

A repeated factor

\[
(s-a)^m
\]

requires every power

\[
\frac1{s-a},
\ldots,
\frac1{(s-a)^m}.
\]

An irreducible quadratic requires a linear numerator.

%%%%%%%%%%%%%%%%%%%%%%%%%%%%%%%%%%%%%%%%%%%%%%%%%%%%%%%%%%%%
\subsubsection{Ignoring Region of Convergence}

A formal expression such as

\[
\frac1{s-a}
\]

does not completely describe bilateral-transform convergence.

The region of convergence carries additional information.

For elementary one-sided Laplace IVPs, the usual convergence half-plane is
often implicit.

%%%%%%%%%%%%%%%%%%%%%%%%%%%%%%%%%%%%%%%%%%%%%%%%%%%%%%%%%%%%
\subsubsection{Using the Final Value Theorem Without Checking Stability}

The equality

\[
\lim_{t\to\infty}f(t)
=
\lim_{s\to0}sF(s)
\]

does not hold for arbitrary transforms.

Oscillatory or unstable systems may violate its hypotheses.

%%%%%%%%%%%%%%%%%%%%%%%%%%%%%%%%%%%%%%%%%%%%%%%%%%%%%%%%%%%%
\subsubsection{Confusing Poles with Zeros}

Poles are associated with denominator roots.

Zeros are associated with numerator roots after relevant cancellation.

They play different roles in system behavior.

%%%%%%%%%%%%%%%%%%%%%%%%%%%%%%%%%%%%%%%%%%%%%%%%%%%%%%%%%%%%
\subsubsection{Assuming Every Pole on the Imaginary Axis Is Stable}

Simple imaginary-axis poles can correspond to sustained oscillation.

Repeated imaginary-axis poles can produce growing terms.

Stability classification requires more than locating a pole exactly on
the axis.

%%%%%%%%%%%%%%%%%%%%%%%%%%%%%%%%%%%%%%%%%%%%%%%%%%%%%%%%%%%%
\subsubsection{Mixing Fourier Transform Conventions}

Different texts distribute factors of

\[
2\pi
\]

differently between forward and inverse transforms.

All formulas must use one convention consistently.

%%%%%%%%%%%%%%%%%%%%%%%%%%%%%%%%%%%%%%%%%%%%%%%%%%%%%%%%%%%%
\subsubsection{Applying Transform Tables Without Algebraic Preparation}

Expressions often require:

\[
\boxed{
\text{factoring},
}
\]

\[
\boxed{
\text{completing the square},
}
\]

or

\[
\boxed{
\text{partial fractions}
}
\]

before a standard transform can be recognized.

%%%%%%%%%%%%%%%%%%%%%%%%%%%%%%%%%%%%%%%%%%%%%%%%%%%%%%%%%%%%
\subsection{Applications}
\label{subsec:transform-method-applications}
%%%%%%%%%%%%%%%%%%%%%%%%%%%%%%%%%%%%%%%%%%%%%%%%%%%%%%%%%%%%

\begin{applicationbox}

\textbf{Differential Equations.}
Laplace transforms convert linear initial-value problems into algebraic
equations.

\medskip

\textbf{Mechanical Systems.}
Step and impulse forcing can represent sudden loads, impacts, and
switching forces.

\medskip

\textbf{Electrical Engineering.}
Circuit equations are analyzed using transfer functions, poles, zeros,
and frequency response.

\medskip

\textbf{Control Theory.}
Transfer functions describe input-output relationships and stability of
linear systems.

\medskip

\textbf{Signal Processing.}
Fourier methods decompose signals into frequencies, while convolution
describes filtering.

\medskip

\textbf{Partial Differential Equations.}
Fourier, sine, cosine, and Laplace transforms reduce PDEs to simpler
ordinary differential or algebraic equations.

\medskip

\textbf{Probability.}
Transforms such as Laplace transforms, characteristic functions, and
moment-generating functions encode probability distributions.

\medskip

\textbf{Discrete Systems.}
The \(z\)-transform converts difference equations into algebraic
relations.

\end{applicationbox}

%%%%%%%%%%%%%%%%%%%%%%%%%%%%%%%%%%%%%%%%%%%%%%%%%%%%%%%%%%%%
\subsection{Exercises}
\label{subsec:transform-method-exercises}
%%%%%%%%%%%%%%%%%%%%%%%%%%%%%%%%%%%%%%%%%%%%%%%%%%%%%%%%%%%%

\begin{exercise}[1]
Define the one-sided Laplace transform.
\end{exercise}

\begin{exercise}[1]
Define exponential order.
\end{exercise}

\begin{exercise}[1]
State the linearity property of the Laplace transform.
\end{exercise}

\begin{exercise}[1]
State the transform of \(1\).
\end{exercise}

\begin{exercise}[1]
State the transform of \(t^n\).
\end{exercise}

\begin{exercise}[1]
State the transform of \(e^{at}\).
\end{exercise}

\begin{exercise}[1]
State the transforms of \(\sin bt\) and \(\cos bt\).
\end{exercise}

\begin{exercise}[1]
State the transform of \(f'(t)\).
\end{exercise}

\begin{exercise}[1]
Define the Heaviside step function.
\end{exercise}

\begin{exercise}[1]
Define convolution.
\end{exercise}

\begin{exercise}[1]
Define a transfer function.
\end{exercise}

\begin{exercise}[1]
Define a pole of a rational transfer function.
\end{exercise}

\begin{exercise}[2]
Find

\[
\mathcal{L}\{4+3t-2t^2\}.
\]
\end{exercise}

\begin{exercise}[2]
Find

\[
\mathcal{L}\{e^{5t}\}.
\]
\end{exercise}

\begin{exercise}[2]
Find

\[
\mathcal{L}\{\cos4t\}.
\]
\end{exercise}

\begin{exercise}[2]
Find

\[
\mathcal{L}\{\sin7t\}.
\]
\end{exercise}

\begin{exercise}[2]
Find

\[
\mathcal{L}\{e^{2t}\cos3t\}.
\]
\end{exercise}

\begin{exercise}[2]
Find

\[
\mathcal{L}^{-1}
\left\{
\frac1{s+4}
\right\}.
\]
\end{exercise}

\begin{exercise}[2]
Find

\[
\mathcal{L}^{-1}
\left\{
\frac5{s^2+25}
\right\}.
\]
\end{exercise}

\begin{exercise}[2]
Find

\[
\mathcal{L}^{-1}
\left\{
\frac{s}{s^2+9}
\right\}.
\]
\end{exercise}

\begin{exercise}[2]
Find the Laplace transform of

\[
u_3(t).
\]
\end{exercise}

\begin{exercise}[2]
Find

\[
\mathcal{L}^{-1}
\left\{
\frac{
e^{-2s}
}{
s
}
\right\}.
\]
\end{exercise}

\begin{exercise}[2]
Find

\[
\mathcal{L}
\{
\delta(t-5)
\}.
\]
\end{exercise}

\begin{exercise}[2]
Find the poles of

\[
H(s)
=
\frac{
s+3
}{
(s+1)(s+4)
}.
\]
\end{exercise}

\begin{exercise}[2]
Find the zero of the previous transfer function.
\end{exercise}

\begin{exercise}[3]
Derive

\[
\mathcal{L}\{t^n\}
=
\frac{n!}{s^{n+1}}.
\]
\end{exercise}

\begin{exercise}[3]
Derive

\[
\mathcal{L}\{f'\}
=
sF-f(0).
\]
\end{exercise}

\begin{exercise}[3]
Derive the second-derivative transform.
\end{exercise}

\begin{exercise}[3]
Derive the general \(n\)-th derivative transform.
\end{exercise}

\begin{exercise}[3]
Derive the first shifting theorem.
\end{exercise}

\begin{exercise}[3]
Derive the second shifting theorem.
\end{exercise}

\begin{exercise}[3]
Solve

\[
y'+3y=6,
\qquad
y(0)=0,
\]

using Laplace transforms.
\end{exercise}

\begin{exercise}[3]
Solve

\[
y''+9y=0,
\qquad
y(0)=2,
\qquad
y'(0)=3,
\]

using Laplace transforms.
\end{exercise}

\begin{exercise}[3]
Solve

\[
y''-y=e^t,
\qquad
y(0)=0,
\qquad
y'(0)=0,
\]

using Laplace transforms.
\end{exercise}

\begin{exercise}[3]
Find

\[
\mathcal{L}^{-1}
\left\{
\frac{
2s+3
}{
(s+1)(s+2)
}
\right\}.
\]
\end{exercise}

\begin{exercise}[3]
Invert

\[
\frac1{
s(s+1)^2
}
\]

using partial fractions.
\end{exercise}

\begin{exercise}[3]
Rewrite the function

\[
f(t)
=
\begin{cases}
2,
&
0\leq t<4,
\\
7,
&
t\geq4
\end{cases}
\]

using a Heaviside step function.
\end{exercise}

\begin{exercise}[3]
Find the Laplace transform of the previous piecewise function.
\end{exercise}

\begin{exercise}[3]
Find

\[
\mathcal{L}^{-1}
\left\{
e^{-3s}
\frac{
s
}{
s^2+4
}
\right\}.
\]
\end{exercise}

\begin{exercise}[3]
Solve an ODE with forcing

\[
u_2(t)
\]

using the second shifting theorem.
\end{exercise}

\begin{exercise}[3]
Solve

\[
y''+4y
=
\delta(t-2),
\]

with zero initial conditions.
\end{exercise}

\begin{exercise}[3]
Prove that convolution is commutative.
\end{exercise}

\begin{exercise}[3]
Use convolution to find

\[
\mathcal{L}^{-1}
\left\{
\frac1{
s(s^2+4)
}
\right\}.
\]
\end{exercise}

\begin{exercise}[3]
Derive the transfer function of

\[
y''+5y'+6y=u(t)
\]

under zero initial conditions.
\end{exercise}

\begin{exercise}[3]
Find the impulse response of the previous system.
\end{exercise}

\begin{exercise}[3]
Find the step response of the previous system.
\end{exercise}

\begin{exercise}[3]
Explain why poles of a constant-coefficient ODE transfer function are
related to characteristic roots.
\end{exercise}

\begin{exercise}[3]
Derive

\[
\mathcal{L}
\{
e^{At}
\}
=
(sI-A)^{-1}.
\]
\end{exercise}

\begin{exercise}[3]
Derive the variation-of-constants formula using Laplace transforms.
\end{exercise}

\begin{exercise}[3]
Explain the conditions needed before using the final value theorem.
\end{exercise}

\begin{exercise}[3]
Derive the Fourier transform derivative property under the convention
used in this section.
\end{exercise}

\begin{exercise}[3]
Explain why convolution becomes multiplication after a Fourier transform.
\end{exercise}

\begin{exercise}[3]
Compare the Laplace, Fourier, and \(z\)-transforms conceptually.
\end{exercise}

\begin{exercise}[4]
Study uniqueness of the inverse Laplace transform.
\end{exercise}

\begin{exercise}[4]
Derive the Bromwich inversion integral.
\end{exercise}

\begin{exercise}[4]
Study regions of convergence for bilateral Laplace transforms.
\end{exercise}

\begin{exercise}[4]
Investigate complex poles and residues for inverse Laplace transforms.
\end{exercise}

\begin{exercise}[4]
Derive inverse transforms using the residue theorem.
\end{exercise}

\begin{exercise}[4]
Study repeated poles and polynomial-exponential responses.
\end{exercise}

\begin{exercise}[4]
Derive the periodic-function Laplace transform formula.
\end{exercise}

\begin{exercise}[4]
Develop distributional derivatives of Heaviside functions and show that

\[
\frac{d}{dt}u_a(t)
=
\delta(t-a)
\]

in the distributional sense.
\end{exercise}

\begin{exercise}[4]
Study Green's functions through impulse responses.
\end{exercise}

\begin{exercise}[4]
Derive transfer matrices for multivariable linear systems.
\end{exercise}

\begin{exercise}[4]
Study controllability and observability through transform-domain
representations.
\end{exercise}

\begin{exercise}[4]
Investigate pole-zero cancellation and its effect on internal stability.
\end{exercise}

\begin{exercise}[4]
Study frequency response by evaluating transfer functions on

\[
s=i\omega.
\]
\end{exercise}

\begin{exercise}[4]
Derive the Fourier inversion formula under suitable hypotheses.
\end{exercise}

\begin{exercise}[4]
Study the Plancherel and Parseval identities.
\end{exercise}

\begin{exercise}[4]
Derive Fourier sine and cosine transform inversion formulas.
\end{exercise}

\begin{exercise}[4]
Use Fourier transforms to solve a differential equation on the real line.
\end{exercise}

\begin{exercise}[4]
Derive the one-sided \(z\)-transform shift formulas.
\end{exercise}

\begin{exercise}[4]
Use the \(z\)-transform to solve a linear recurrence relation.
\end{exercise}

%%%%%%%%%%%%%%%%%%%%%%%%%%%%%%%%%%%%%%%%%%%%%%%%%%%%%%%%%%%%
\subsection{Conceptual Summary}
\label{subsec:transform-method-summary}
%%%%%%%%%%%%%%%%%%%%%%%%%%%%%%%%%%%%%%%%%%%%%%%%%%%%%%%%%%%%

The Laplace transform is

\[
\boxed{
\mathcal{L}\{f(t)\}
=
\int_0^\infty
e^{-st}f(t)\,dt.
}
\]

It converts derivatives into algebraic expressions:

\[
\boxed{
\mathcal{L}\{f'\}
=
sF-f(0),
}
\]

and

\[
\boxed{
\mathcal{L}\{f''\}
=
s^2F
-
sf(0)
-
f'(0).
}
\]

The first shifting theorem is

\[
\boxed{
\mathcal{L}
\{
e^{at}f(t)
\}
=
F(s-a).
}
\]

The second shifting theorem is

\[
\boxed{
\mathcal{L}
\{
u_a(t)f(t-a)
\}
=
e^{-as}F(s).
}
\]

Convolution is

\[
\boxed{
(f*g)(t)
=
\int_0^t
f(\tau)g(t-\tau)\,d\tau,
}
\]

and satisfies

\[
\boxed{
\mathcal{L}\{f*g\}
=
FG.
}
\]

For zero initial conditions, the transfer function is

\[
\boxed{
H(s)
=
\frac{
Y(s)
}{
U(s)
}.
}
\]

The Fourier transform under the convention used here is

\[
\boxed{
\widehat{f}(\omega)
=
\int_{-\infty}^{\infty}
f(t)e^{-i\omega t}\,dt.
}
\]

Transform methods therefore convert

\[
\boxed{
\text{differentiation}
\rightarrow
\text{multiplication},
}
\]

\[
\boxed{
\text{convolution}
\rightarrow
\text{multiplication},
}
\]

and

\[
\boxed{
\text{time delay}
\rightarrow
\text{exponential factors}.
}
\]

%%%%%%%%%%%%%%%%%%%%%%%%%%%%%%%%%%%%%%%%%%%%%%%%%%%%%%%%%%%%
\subsection{Section Connections}
\label{subsec:transform-method-connections}
%%%%%%%%%%%%%%%%%%%%%%%%%%%%%%%%%%%%%%%%%%%%%%%%%%%%%%%%%%%%

\chapterconnection{
Transform Methods extends the ordinary differential-equation techniques
of Section~\ref{sec:ordinary-differential-equations} by converting
differential equations into algebraic equations in a transform variable.
Laplace transforms are particularly effective for initial-value problems,
piecewise forcing, impulses, linear systems, and transfer functions.
Convolution connects transform theory with Green's functions and
input-output systems, while Fourier transforms lead naturally to
frequency-domain analysis and partial differential equations. The
\(z\)-transform provides an analogous framework for discrete dynamical
systems and difference equations. The next section, Dynamical Systems,
moves from explicit solution techniques to qualitative analysis of
equilibria, phase portraits, stability, bifurcations, and long-term
behavior.
}

%%%%%%%%%%%%%%%%%%%%%%%%%%%%%%%%%%%%%%%%%%%%%%%%%%%%%%%%%%%%
% SECTION SOURCE TRACKING
%%%%%%%%%%%%%%%%%%%%%%%%%%%%%%%%%%%%%%%%%%%%%%%%%%%%%%%%%%%%

%------------------------------------------------------------
% 9.2 Transform Methods
%------------------------------------------------------------
%
% Source basis:
%   - Standard Laplace-transform theory.
%   - Standard Fourier-transform theory.
%   - Standard linear systems and operational methods.
%
% Used for:
%   - integral transforms;
%   - Laplace transform definition;
%   - existence and exponential order;
%   - transform linearity;
%   - elementary transform tables;
%   - inverse Laplace transforms;
%   - derivative transforms;
%   - integral transforms;
%   - differentiation with respect to s;
%   - first shifting theorem;
%   - partial fractions;
%   - initial-value problems;
%   - Heaviside step functions;
%   - second shifting theorem;
%   - delayed forcing;
%   - periodic functions;
%   - Dirac delta distributions;
%   - impulsive forcing;
%   - convolution;
%   - convolution theorem;
%   - impulse and step responses;
%   - transfer functions;
%   - poles and zeros;
%   - systems of differential equations;
%   - resolvent matrices;
%   - matrix exponentials;
%   - variation of constants;
%   - initial and final value theorems;
%   - bilateral Laplace transform preview;
%   - Fourier transforms;
%   - derivative and convolution properties;
%   - sine and cosine transforms;
%   - frequency-domain interpretation;
%   - z-transform preview;
%   - applications and exercises.
%
% Notes:
%   - Dynamical Systems is developed next in Section 9.3.
%   - Partial Differential Equations later use Fourier and transform
%     methods in greater depth.
%   - Difference Equations later develop the z-transform connection.
%   - Numerical transform methods are developed more deeply in
%     Chapter 11.
%
%%%%%%%%%%%%%%%%%%%%%%%%%%%%%%%%%%%%%%%%%%%%%%%%%%%%%%%%%%%%